\chapter{基于汇编的裸机代码}
\label{chap:boot}
\setcounter{secnumdepth}{3}

\begin{introduction}[本章内容]
  \item 上电、链接地址与装载地址
  \item 异常级别：从 ARM9 的七种模式到 AArch64 的 EL0～EL3
  \item 不依赖任何环境的早期串口
  \item 异常和中断：向量表、现场保存、中断控制器与定时器
\end{introduction}

\ref{sec:xv6-startup} 节讲本项目的启动代码：从固件交出控制权、统一异常级别，到点亮早期串口。\ref{sec:irq} 节讲异常和中断：向量表、现场保存、两级中断控制器和通用定时器。清 BSS、建启动页表、打开 MMU 和多核启动与页表关系密切，放在第~\ref{chap:mmu} 章的 \ref{sec:bootpgt}、\ref{sec:smp-boot} 节。这两节的代码全部来自 xv6-aarch64；它要兼容真机固件、QEMU 和多核，比教学内核复杂，所以每一步都先讲清楚背后的原理，这些原理主要参考了 Sergey Matyukevich 的 \emph{raspberry-pi-os}\cite{rpios}（下文简称 RPi OS）第 1～3 课。RPi OS 后三课的原理也按同样的方式融入了后面各章：第 4 课（进程调度器）在第~\ref{chap:process} 章，第 5 课（用户进程与系统调用）在第~\ref{chap:user} 章，第 6 课（虚拟内存管理）在第~\ref{chap:mmu} 章，那些地方同样只用本项目的代码来讲。

\section{xv6-aarch64 的启动代码}
\label{sec:xv6-startup}
% 3.1 xv6-aarch64 的启动代码（融入 raspberry-pi-os 第 1、2 课的原理）
\subsection{上电之后发生了什么}

树莓派的启动与 PC 不同：最先运行的不是 ARM 核，而是 VideoCore GPU。

\begin{enumerate}
  \item GPU 内部 ROM 从 SD 卡的 FAT32 分区读取 \code{bootcode.bin}，再加载 \code{start.elf}（GPU 固件）；
  \item \code{start.elf} 读取 \code{config.txt}，初始化 SDRAM 和时钟，把内核镜像 \code{kernel8-xv6\_wifi.img} 拷到物理地址 \code{0x80000}；
  \item 固件还在物理地址 0 处放一段很小的 \emph{armstub}。它让四个 Cortex-A53 核一起开始执行：CPU0 跳到 \code{0x80000}，CPU1～3 在 \emph{spin-table} 上循环等待——各自读取 \code{0xe0}、\code{0xe8}、\code{0xf0} 处的 64 位地址，非零就跳过去；
  \item 固件自带的 64 位 armstub 会先把核从 EL3 降到 EL2，因此内核通常在 EL2 拿到控制权；QEMU 也是如此。
\end{enumerate}

项目早期曾使用自制的极简 armstub（\file{kernel/armstub.S}），它只做一次跳转、保留在 EL3。后来发现这样做会让固件的 property mailbox 在真机上完全没有响应——而 SD 时钟查询、USB 上电、摄像头供电都依赖 mailbox。所以现在直接使用固件自带的 armstub，内核启动代码则要能同时处理 EL3、EL2 两种入口。

\paragraph{固件、\code{config.txt} 与镜像文件名.} 树莓派的固件是闭源的，ARM 一侧能控制它的只有启动分区里的 \code{config.txt}。从一张空白 SD 卡启动一个裸机内核，至少需要四个文件：GPU 引导程序 \code{bootcode.bin}、GPU 固件 \code{start.elf}、\code{config.txt} 和内核镜像。固件根据镜像文件名决定 CPU 的执行状态：以 \code{8} 结尾的 \code{kernel8.img} 表示 ARMv8，固件会让四个核以 AArch64 状态启动；本项目用 \code{kernel=} 指定了自己的文件名，同时写了 \code{arm\_64bit=1}。

\paragraph{另一种启动方式：\code{kernel\_old}.} 最简单的裸机教程（如 raspberry-pi-os\cite{rpios}）在 \code{config.txt} 中写 \code{kernel\_old=1}。这时固件把镜像装在物理地址 0，并且不安装 armstub：四个核都从地址 0 开始执行同一段内核代码，而且停留在 EL3。内核必须自己读 \reg{MPIDR\_EL1} 的低 8 位区分核号，把非 0 号核挂进死循环，再自己从 EL3 降到 EL1。这种方式的好处是“什么都由自己掌控”，代价是失去了固件 armstub 提供的 spin-table 多核唤醒协议，也失去了 armstub 对 EL3 安全状态的初始化。本项目选择相反的路线：让固件做它擅长的事，内核只在 \code{\_entry} 中兼容两种入口级别。

\subsection{链接地址与装载地址}

内核在高地址 \code{0xffffff8000080000} 链接，却被装到物理地址 \code{0x80000}。\file{kernel/memlayout.h} 用两个宏描述这种对应关系：

\begin{ccode}
#define KERNBASE  0xffffff8000000000L     // 内核虚拟地址空间起点
#define KERNLINK  (KERNBASE + KERNPA)     // 内核链接地址
#define V2P(a) (((uint64)(a)) - KERNBASE)
#define P2V(a) ((void *)(((char *)(a)) + KERNBASE))
\end{ccode}

\code{KERNBASE} 选在 39 位地址空间的高半部分起点：当 \reg{TCR\_EL1.T1SZ}$=25$ 时，\reg{TTBR1\_EL1} 负责的范围恰好是 \code{0xffffff8000000000}～\code{0xffffffffffffffff}。内核把全部物理内存线性映射到 \code{KERNBASE + PA}，所以物理地址与内核虚拟地址之间只差一个常数。

MMU 打开之前，CPU 用物理地址取指。此时如果执行 \code{ldr x0, =symbol} 得到的是高虚拟地址，访问它会立即出错；只有 PC 相对寻址的 \code{adr}/\code{adrp} 能得到正确的物理地址。这一限制决定了 \file{entry.S} 前半部分的写法。

\paragraph{独立环境下的编译.} 内核不能依赖任何运行时环境，编译时要明确告诉编译器这一点：不链接 C 标准库（库中的大多数函数最终都要调用操作系统，而这里下面已经没有操作系统了）；不使用标准启动文件（设置栈、初始化静态数据、调用 \code{main} 这些事都由启动汇编自己完成）；声明为独立（freestanding）环境，使编译器不对 \code{memcpy} 之类的函数做任何假设；只使用通用寄存器，不让编译器生成 NEON/浮点指令，否则上下文切换就必须额外保存这些寄存器。本项目对应的选项是 \code{-ffreestanding -nostdlib -fno-common} 和 \code{-mcpu=cortex-a53+nofp}（第~\ref{chap:organization} 章）。

\paragraph{为什么要把 ELF 变成裸镜像.} 链接器输出的是 ELF 文件，它需要一个“装载器”去解析段表、把各段放到指定地址。固件不做这件事，它只是把文件按字节原样拷到内存，然后从第一个字节开始执行。所以要用 \code{objcopy -O binary} 抽出所有代码段和数据段，拼成一个裸镜像，并且保证启动代码位于镜像的最开头：常见的做法是把启动代码放进单独的 \code{.text.boot} 段并在链接脚本中排在第一位；本项目则让 \file{entry.o} 成为 \code{OBJS} 的第一个目标文件，\code{\_entry} 自然落在 \code{.text} 的起点。

\subsection{统一异常级别}
\label{sec:xv6-el}

AArch64 有四个异常级别：EL0 运行用户程序，EL1 运行操作系统内核，EL2 运行虚拟机监控器，EL3 运行安全监控固件。xv6 只需要 EL0 和 EL1。

\paragraph{为什么需要异常级别.} 异常级别可以看成处理器的几种执行模式，每种模式只能使用一部分指令和寄存器。操作系统要实现进程隔离，就必须让用户程序运行在特权最低的 EL0：在这一级只能使用通用寄存器、栈指针和普通的访存指令，不能修改页表基址等系统寄存器，因此一个进程既看不到也改不了别的进程的内存。操作系统本身运行在 EL1，可以配置虚拟内存和大多数系统寄存器。EL2 留给虚拟机监控器：宿主系统在 EL2 隔离各个运行在 EL1 的客户系统，就像操作系统隔离用户进程一样。EL3 负责在“安全世界”和“非安全世界”之间切换，两边的隔离由硬件保证，非安全世界的软件无论如何都访问不到安全世界的指令和数据。

\paragraph{ARM9 时代：七种处理器模式.} 四个异常级别是 ARMv8 才有的说法。在 ARM9（ARM920T、ARM926EJ-S 等 ARMv4T/v5 内核，S3C2440、AT91 这类板子上最常见）时代，处理器用的是\emph{处理器模式}（processor mode），当前模式记在 \reg{CPSR} 的低 5 位 M[4:0] 里，一共七种（表~\ref{tab:arm-modes}）：用户模式（User）运行应用程序；系统模式（System）是“有特权的用户模式”；其余五种各对应一类异常——复位和 \code{swi} 指令进入管理模式（Supervisor，SVC），未定义指令进入未定义模式（Undefined），取指或数据访问出错进入中止模式（Abort），普通中断和快速中断分别进入 IRQ 模式和 FIQ 模式。

\paragraph{模式不等于特权级.} 七种模式里只有 User 是非特权的，其余六种的特权完全相同，都能改 \reg{CPSR}、访问 CP15 协处理器（MMU、缓存的控制寄存器都在这里）。所以 ARM9 实际上只有两个特权级，模式之间的区别不在于“能做什么”，而在于\emph{用哪一组寄存器}：每个异常模式有自己私有（banked）的 \code{r13}（SP）、\code{r14}（LR）和 \reg{SPSR}，FIQ 还额外私有 \code{r8}～\code{r12}，这样快速中断处理程序不用保存任何通用寄存器就能开工。异常的类型直接决定进入哪个模式；向量表固定在地址 \code{0x00000000}（或 CP15 选择的高端 \code{0xffff0000}），每类异常占一个 4 字节的槽，只放得下一条跳转指令。

这种设计给操作系统带来不少麻烦。内核必须在启动时依次切进 IRQ、Abort、Undefined 各模式，为每个模式单独设置栈指针；而内核真正想在一个统一的上下文里处理所有陷入，所以 Linux 的 ARM 移植在每个异常入口只往该模式的小栈里存三个字（\code{r0}、\code{lr}、\code{spsr}），随即切换到 SVC 模式，在进程的内核栈上继续处理——七种模式在软件里又被折回成“用户态”和“内核态”两种。另外，特权代码可以直接用 \code{msr cpsr\_c, \#0xd3} 改写 \reg{CPSR} 进入任意模式，模式切换并不一定经过异常。

\paragraph{TrustZone 与虚拟化：再加两种模式.} ARMv6K 和 ARMv7 引入安全扩展（TrustZone），处理器多了一个“安全/非安全”状态位，以及负责在两个世界之间切换的监控模式（Monitor），由 \code{smc} 指令进入。ARMv7 的虚拟化扩展（Cortex-A15 起）又加了一个 Hyp 模式，供虚拟机监控器使用，由 \code{hvc} 进入。为了说清楚这些模式之间的高低，ARMv7 手册开始使用\emph{特权级}（PL0：User；PL1：六种原有特权模式加 Monitor；PL2：Hyp）。到这时已经有九种模式，而“特权高低”和“用哪组寄存器”这两件事仍然搅在一起。

\paragraph{AArch64：把模式折叠成异常级别.} ARMv8 的 AArch64 状态把这两件事拆开了。\emph{异常级别}只表示特权的高低：EL0 对应原来的 User，EL1 一次吸收了 SVC、IRQ、FIQ、Abort、Undefined、System 六种模式，EL2 对应 Hyp，EL3 对应 Monitor。异常的\emph{类型}不再决定模式，而由两样东西表达：

\begin{itemize}
  \item 异常发生时跳到目标级别向量表里的哪一项。\reg{VBAR\_ELx} 指向的表有 16 项，按“来自哪一级、用哪个栈 $\times$ 同步/IRQ/FIQ/SError”排列，每项 128 字节，能直接放下一段保存现场的代码（\ref{sec:irq} 节）；
  \item 同步异常的具体原因写进 \reg{ESR\_ELx} 的 EC 字段：\code{svc} 是 \code{0x15}，来自 EL0 的数据中止是 \code{0x24}，未定义指令是 \code{0x00}……原来要靠“当前处于 Abort 模式还是 Undefined 模式”来区分的事情，现在读一个寄存器就知道。
\end{itemize}

私有寄存器也大大精简：31 个通用寄存器不再分组，每一级只保留自己的栈指针 \reg{SP\_ELx}、返回地址 \reg{ELR\_ELx} 和保存状态 \reg{SPSR\_ELx}，FIQ 的快速寄存器组被取消。原来的 System 模式（特权级、但用 User 的寄存器）也有对应物：EL1t 在 EL1 下借用 \reg{SP\_EL0}，EL1h 使用自己的 \reg{SP\_EL1}（见下文 \code{SPSR} 的模式字段）。最重要的一点是：AArch64 里没有任何指令能直接改写当前级别，升级只能靠异常，降级只能靠 \code{eret}。

\begin{table}[htbp]
  \centering
  \caption{ARM9/ARMv7 处理器模式与 AArch64 异常级别}
  \label{tab:arm-modes}
  \small
  \begin{tabularx}{\linewidth}{@{}llllX@{}}
    \toprule
    模式 & \reg{CPSR}.M & 出现于 & AArch64 & 何时进入（旧架构） \\
    \midrule
    User (usr)       & \code{10000} & ARMv4 & EL0 & 应用程序运行 \\
    FIQ (fiq)        & \code{10001} & ARMv4 & EL1 & 快速中断；私有 \code{r8}～\code{r14} \\
    IRQ (irq)        & \code{10010} & ARMv4 & EL1 & 普通中断 \\
    Supervisor (svc) & \code{10011} & ARMv4 & EL1 & 复位、\code{swi}/\code{svc} 系统调用 \\
    Abort (abt)      & \code{10111} & ARMv4 & EL1 & 预取中止、数据中止 \\
    Undefined (und)  & \code{11011} & ARMv4 & EL1 & 未定义指令 \\
    System (sys)     & \code{11111} & ARMv4 & EL1 & 无对应异常，由特权代码直接切入 \\
    Monitor (mon)    & \code{10110} & ARMv6K/v7 & EL3 & \code{smc}，安全/非安全世界切换 \\
    Hyp (hyp)        & \code{11010} & ARMv7VE & EL2 & \code{hvc}、虚拟化陷入 \\
    \bottomrule
  \end{tabularx}
\end{table}

这张对应关系并不只是类比。ARMv8 处理器仍能在 EL0、EL1 上运行 32 位代码（AArch32 状态），这时旧的模式依旧存在，而且严格按表~\ref{tab:arm-modes} 落在对应的异常级别上；那些私有寄存器也被映射到 64 位通用寄存器上（例如 IRQ 模式的 \code{r13}/\code{r14} 就是 \code{X17}/\code{X16}），所以 64 位内核切换到 32 位进程时不需要另外保存它们。对比早期面向树莓派 1（ARMv6）的 xv6 移植：它要在启动时为 IRQ、Abort、Undefined 各设一个栈，并在每个入口切回 SVC 模式；本项目在 AArch64 上只需要每个核一个 \reg{SP\_EL1}，所有异常都在同一个 EL1 里处理。

\paragraph{异常级别只能靠异常来改变.} 没有更高特权软件的参与，程序无法提升自己的级别——否则隔离就形同虚设。级别只在异常发生和返回时改变。发生异常（非法访存、\code{svc} 指令、硬件中断等）并由 EL$n$ 处理时，硬件依次做三件事：把返回地址存入 \reg{ELR\_ELn}，把当时的处理器状态存入 \reg{SPSR\_ELn}，跳到异常向量；处理程序结束时执行 \code{eret}，硬件再从这两个寄存器恢复状态和地址。关键在于这两个寄存器都是可写的，处理程序不一定要“回到原处”。启动时降级正是利用这一点。

调试启动代码时，一个简单而有用的手段是读出 \reg{CurrentEL}：当前级别保存在它的第 [3:2] 位，右移 2 位即得 0～3。用不同的固件或 \code{config.txt} 启动时，读到的值可能是 3 或 2，下面的代码对两者都要处理：

\file{kernel/entry.S} 的第一件事就是读 \reg{CurrentEL}，不论从 EL3 还是 EL2 进入，都降到 EL1：

\begin{asmcode}
_entry:
        mrs x15, mpidr_el1
        and x15, x15, #0xff          // 记下核号，后面区分主核与副核
        mrs x0, CurrentEL
        cmp x0, #0xc                 // CurrentEL[3:2] == 3 ?
        b.ne 2f
        mov x0, #(1 << 31)           // HCR_EL2.RW：EL1 使用 AArch64
        msr hcr_el2, x0
        msr cntvoff_el2, xzr         // 虚拟计数器偏移清零
        mov x0, #3
        msr cnthctl_el2, x0          // 允许 EL1 访问物理计数器/定时器
        mov x0, #0x431               // SCR_EL3：RES1 | RW | NS
        msr scr_el3, x0
        mov x0, #0x3c5               // 目标 EL1h，屏蔽 D/A/I/F
        msr spsr_el3, x0
        adr x0, 3f
        msr elr_el3, x0
        eret
2:      cmp x0, #0x8                 // EL2 路径与上面相同，只是写 *_el2
        ...
        eret
3:      ldr x0, =((3 << 28) | (3 << 22) | (1 << 20) | (1 << 11))
        msr sctlr_el1, x0            // 只保留 RES1 位：MMU、缓存先全部关闭
        isb
        cbnz x15, secondary_wait     // 副核去等待
        bl early_mini_uart           // 主核点亮串口
        b entry
\end{asmcode}

“降级”的方式很有意思：处理器没有一条“切换到 EL1”的指令，只能伪造一次异常返回。把目标状态写进 \reg{SPSR\_ELx}、把返回地址写进 \reg{ELR\_ELx}，然后执行 \code{eret}，CPU 就会“返回”到 EL1 的标号 \code{3} 处。表~\ref{tab:el-regs} 列出了涉及的寄存器。

\begin{table}[htbp]
  \centering
  \caption{降级到 EL1 时设置的系统寄存器}
  \label{tab:el-regs}
  \small
  \begin{tabularx}{\linewidth}{@{}llX@{}}
    \toprule
    寄存器 & 值 & 含义 \\
    \midrule
    \reg{HCR\_EL2} & bit 31 & RW=1：下一级 EL1 以 AArch64 运行，否则会掉进 AArch32 \\
    \reg{CNTVOFF\_EL2} & 0 & 虚拟计数器 = 物理计数器，内核用的是虚拟定时器 \\
    \reg{CNTHCTL\_EL2} & 3 & EL1/EL0 可以访问物理计数器和物理定时器 \\
    \reg{SCR\_EL3} & \code{0x431} & 非安全状态，下一级为 AArch64 \\
    \reg{SPSR\_ELx} & \code{0x3c5} & M[3:0]=\code{0101} 即 EL1h（使用 \reg{SP\_EL1}）；DAIF 四位全部置 1 \\
    \reg{SCTLR\_EL1} & RES1 位 & 不继承固件留下的 MMU/缓存状态 \\
    \bottomrule
  \end{tabularx}
\end{table}

\code{SPSR} 中的 D、A、I、F 四位分别屏蔽调试异常、SError、IRQ 和 FIQ。启动期间必须全部屏蔽：此时既没有异常向量表，也没有栈，任何一个中断都会让 CPU 跳到不可预知的地方。

\paragraph{\code{SPSR} 里保存的处理器状态.} 写入 \reg{SPSR\_ELx} 的值就是 \code{eret} 之后的 PSTATE，它由三部分组成：条件标志 N、Z、C、V（上一次运算结果是否为负、为零、无符号溢出、有符号溢出，条件分支指令据此跳转）；中断屏蔽位 D、A、I、F；以及模式字段 M[3:0]，它决定返回到哪一级、使用哪个栈指针。EL1 有两种栈模式：EL1h 使用 EL1 自己的 \reg{SP\_EL1}，EL1t 借用 EL0 的 \reg{SP\_EL0}。内核总是使用 EL1h，这样陷入内核时自动换到内核栈，用户栈指针原样留在 \reg{SP\_EL0} 中。有的教学内核只屏蔽 A、I、F 三位（值 \code{7<<6}），本项目连同调试异常 D 一起屏蔽（\code{0xf<<6}）。

\paragraph{\code{SCTLR\_EL1} 的初值.} \reg{SCTLR\_EL1} 控制 EL1 下的 MMU、数据缓存、指令缓存和字节序。手册把其中若干位标为 RES1（保留，必须写 1），即上面那个 \code{(3<<28)|(3<<22)|(1<<20)|(1<<11)}；其余位全写 0，含义是：EL1 和 EL0 的数据访问都是小端，指令缓存、数据缓存和 MMU 全部关闭。不继承固件留下的任何状态，后面的代码才有确定的起点。

\paragraph{\code{HCR\_EL2} 与 \code{SCR\_EL3}.} 即使不实现虚拟机监控器，也必须设置 \reg{HCR\_EL2} 的 RW 位，因为它决定 EL1 运行在 AArch64 还是 AArch32；同理，\reg{SCR\_EL3} 的 RW 位决定 EL2 的执行状态，NS 位让所有更低的级别处于非安全状态。漏掉 RW 位，\code{eret} 之后处理器会以 32 位状态执行 64 位指令，表现为“毫无输出地死掉”。

\subsection{不依赖任何环境的早期串口}
\label{sec:early-uart}

普通的 \code{printf} 依赖一长串设施：有效的栈、清零的 BSS、页表、内存分配器、自旋锁和控制台驱动。如果内核在这些设施就绪之前挂掉，串口上什么也看不到，连“死在哪一步”都无从判断。项目在真机上调通的第一个关键是：用纯汇编写一个不需要栈、不需要内存、不需要 C 运行环境的串口输出。

\paragraph{内存映射的设备寄存器.} BCM2837 上所有外设都通过内存映射寄存器访问：树莓派 3 把 \code{0x3F000000} 以上的一段物理地址留给外设，每个寄存器是 32 位，读写某一位就是配置或启动设备的某项功能。Broadcom 手册中的地址以 \code{0x7E} 开头，那是 VideoCore 总线上的地址，ARM 一侧对应的物理地址要换成 \code{0x3F} 开头（第~\ref{chap:mmu} 章）。

树莓派 3 的 GPIO14/15 可以连到 PL011 UART0 或 AUX Mini UART。本内核使用 Mini UART，即 GPIO14/15 的 ALT5 功能。树莓派有两个 UART：PL011 功能完整，但在树莓派 3 上默认被板载蓝牙占用；Mini UART 结构简单，只需几个寄存器就能工作，代价是它的波特率依赖 VPU 核心时钟。初始化分三步：

\paragraph{第一步：选择 GPIO 复用功能.} 大多数 GPIO 引脚可以连接不同的外设，每个引脚有 0～5 共六种“复用功能”，用 3 位编码，十个引脚共用一个 \code{GPFSELn} 寄存器。GPIO14、15 由 \code{GPFSEL1} 的第 14:12 位和 17:15 位控制；它们的 ALT5 分别是 Mini UART 的 TXD1 和 RXD1，ALT5 的编码是 \code{010}。所以要先把这两个 3 位字段清零，再各写入 2。注意 GPIO 编号与排针针脚号不同：GPIO14 是第 8 脚，GPIO15 是第 10 脚，第 6 脚是地。

\paragraph{第二步：关闭内部上下拉.} 一个悬空的输入引脚读到的值是随机的；上拉或下拉电阻让它在悬空时稳定地读到 1 或 0。串口两根线始终接着对端，不需要上下拉，而且上下拉状态在重启后仍然保持，所以使用前要显式撤销。BCM2837 改变上下拉状态需要一个带时钟的握手过程：先在 \code{GPPUD} 中写入目标状态（0 表示都不要），等待约 150 个周期让控制信号建立；再在 \code{GPPUDCLK0} 中把要修改的引脚对应位写 1，把控制信号“打”进这些引脚，再等 150 个周期保持；最后清除 \code{GPPUD} 和 \code{GPPUDCLK0}。只有收到时钟的引脚会改变，其余引脚保持原状。

\paragraph{第三步：配置 Mini UART.} 顺序很重要：先在 AUX 使能寄存器中打开 Mini UART，因为在此之前它的其他寄存器都不可访问；配置期间先关闭收发器，并永久关闭自动流控（它需要额外的 RTS/CTS 引脚，普通 USB-TTL 线不支持）；关闭收发中断（早期阶段没有中断处理）；把数据位设为 8 位——手册写的是第 0 位，实际硬件要第 0、1 两位都置 1，即写 3；RTS 置为常高；写波特率除数；最后打开收发器。

\code{early\_mini\_uart} 依次完成：

\begin{asmcode}
early_mini_uart:
        mov x13, x30                 // 没有栈，用寄存器保存返回地址
        ldr x9, =0x3f200004          // GPFSEL1：GPIO14/15 = ALT5（功能码 2）
        ldr w10, [x9]
        ldr w11, =0x0003f000
        bic w10, w10, w11
        ldr w11, =0x00012000
        orr w10, w10, w11
        str w10, [x9]
        ...                          // GPPUD/GPPUDCLK0 时序：关闭内部上下拉
        ldr x9, =0x3f215000          // AUX 基址
        ...                          // 使能 AUX、8N1、清 FIFO
        mov w10, #270
        str w10, [x9, #0x68]         // AUX_MU_BAUD：115200 @ 250 MHz
        mov w10, #3
        str w10, [x9, #0x60]         // 打开收发
        mov w12, #'E'
        bl early_mini_uart_putc
        ...
\end{asmcode}

这里用的都是物理地址，因为 MMU 还没打开。Mini UART 的波特率公式是
\[
  \text{baud} = \frac{f_\text{core}}{8\,(\text{divisor}+1)}, \qquad
  \text{divisor} = \frac{250\,000\,000}{8 \times 115\,200} - 1 \approx 270.
\]
它依赖 VPU 核心频率，所以 \file{config.txt} 必须固定 \code{core\_freq=250}。

收发都用轮询：\code{AUX\_MU\_LSR} 的第 5 位为 1 表示发送器空闲，可以向 \code{AUX\_MU\_IO} 写一个字符；第 0 位为 1 表示收到数据，可以从 \code{AUX\_MU\_IO} 读出。早期汇编版只需要发送，所以 \code{early\_mini\_uart\_putc} 只有三条指令：读 LSR、测第 5 位、写 IO。

在没有调试器的裸机上，串口输出几乎是了解程序内部状态的唯一途径。所以一个裸机项目通常要做的第二件事，就是在串口之上移植一个小型 \code{printf}——只需要提供一个发送单个字符的函数。本项目把这件事分成两级：汇编阶段只能输出单个标记字符，C 环境就绪后才由控制台驱动提供完整的 \code{printf}。

此后，启动代码在每个关键节点输出一个字符，形成一串“启动指纹”（表~\ref{tab:boot-marks}）。真机串口上看到的最后一个字符，就指明了内核停在哪一步。

\begin{table}[htbp]
  \centering
  \caption{启动标记字符}
  \label{tab:boot-marks}
  \small
  \begin{tabular}{@{}cl@{\hspace{2em}}cl@{}}
    \toprule
    标记 & 含义 & 标记 & 含义 \\
    \midrule
    \code{E} & 进入 EL1，物理地址串口可用 & \code{1} & \code{kinit1()} 完成 \\
    \code{B} & BSS 清零完成 & \code{2} & 最终内核页表建立 \\
    \code{P} & 启动页表与系统寄存器就绪 & \code{3} & 切换到最终页表 \\
    \code{H} & MMU 已开，跳到高地址并建好栈 & \code{4} & \code{kinit2()} 完成 \\
    \code{0} & 进入 \code{main()} & \code{a}～\code{f} & \code{kalloc.c} 内部细分 \\
    \bottomrule
  \end{tabular}
\end{table}

正常启动时串口先打出 \code{E}，然后是 \code{BPH0abdefc1234}，接着才是 \code{printf} 的第一行 “xv6 kernel is booting”。第~\ref{chap:mmu} 章会讲到一次停在 \code{BPH0abde} 的故障：它说明第一次 \code{kfree()} 已经完成 \code{memset}，却没能拿到 \code{kmem.lock}。

MMU 打开之后，物理地址 \code{0x3f215000} 不再可用。\code{boot\_uart\_mark()} 改用高地址 \code{KERNBASE + 0x3f215000}，所以启动页表里专门为 Mini UART 所在的 2\,MiB 块加了一个设备映射。


% 3.2 异常和中断（全部以 xv6-aarch64 源码讲解）
\section{异常与中断}
\label{sec:irq}

与设备打交道的大多数场合都是异步的：内核向设备发出一条命令，设备并不立即完成，而是在工作结束后“通知”处理器。这种通知打断了正在执行的指令流，迫使处理器转去执行一段处理程序，所以叫“中断”。对操作系统来说，最重要的中断源是定时器：调度器靠周期性的定时器中断度量每个进程运行了多久，并在时间片用完时把 CPU 交给别人。

本节沿着一次中断从发生到返回的路径，依次讲解 xv6-aarch64 的实现：\file{kernel/trapasm.S}（向量表与现场保存）、\file{kernel/trap.c}（分发）、\file{kernel/bcm2837.c}（中断控制器）、\file{kernel/timer.c}（定时器）和 \file{kernel/spinlock.c}（中断屏蔽的嵌套）。

\subsection{异常与中断区别}

在 ARMv8 中，中断只是“异常”的一种。异常分为四类，表~\ref{tab:exc-kinds} 列出了它们在 xv6-aarch64 中的用途。

\begin{table}[htbp]
  \centering
  \caption{四类异常在 xv6-aarch64 中的用途}
  \label{tab:exc-kinds}
  \small
  \begin{tabularx}{\linewidth}{@{}llX@{}}
    \toprule
    类型 & 性质 & 在 xv6-aarch64 中 \\
    \midrule
    同步异常 & 由当前指令引起 & 用户态 \code{svc \#0} 即系统调用；用户态访存错误时打印诊断并杀死进程；内核态同步异常直接 \code{panic} \\
    IRQ & 外部硬件产生，与当前指令无关 & 通用定时器、Mini UART、DWC2 USB、Arasan SDIO、Unicam CSI、核间中断（IPI） \\
    FIQ & 高优先级的“快速中断” & 不使用 \\
    SError & 外部产生的系统错误 & 不处理，进入即停机 \\
    \bottomrule
  \end{tabularx}
\end{table}

同步异常的具体原因记录在 \reg{ESR\_EL1}（异常综合征寄存器）中：高 6 位 [31:26] 是异常类别 EC，低 25 位 ISS 是该类别下的细节；出错的指令地址在 \reg{ELR\_EL1}，出错的数据地址在 \reg{FAR\_EL1}。xv6 只识别一种“正常”的同步异常——EC $=$ \code{0x15}（21），即来自 AArch64 EL0 的 \code{svc}；其余同步异常都按错误处理。内核态的同步异常不尝试恢复：

\begin{ccode}
void kerneltrap()
{
  uint64 esr = r_esr_el1();
  w_esr_el1(0);
  if(intr_get() != 0)
    panic("kerneltrap: interrupts enabled");
  uint64 ec = (esr >> 26) & 0x3f;
  uint64 iss = esr & 0x1ffffff;
  printf("esr %p %p %p\n", esr, ec, iss);
  printf("elr=%p far=%p\n", r_elr_el1(), r_far_el1());
  panic("kerneltrap");
}
\end{ccode}

\subsection{异常向量}

每类异常都要有自己的处理入口；而且同一类异常，从不同的执行状态进入时也要用不同的入口。以运行在 EL1 的内核为视角，进入异常时的执行状态有四种：EL1t（EL1 使用 \reg{SP\_EL0}）、EL1h（EL1 使用自己的 \reg{SP\_EL1}，内核正是这种模式）、来自 64 位 EL0、来自 32 位 EL0。四类异常乘以四种状态，共 16 个入口，它们排成一张\emph{异常向量表}：表必须 2\,KiB 对齐，每个入口占 \code{0x80} 字节（表~\ref{tab:vectors}）。

\begin{table}[htbp]
  \centering
  \caption{\reg{VBAR\_EL1} 向量表布局}
  \label{tab:vectors}
  \small
  \begin{tabular}{@{}lllll@{}}
    \toprule
    偏移 & 来源 & 同步 & IRQ & 本内核 \\
    \midrule
    \code{0x000} & 当前 EL，使用 \reg{SP\_EL0} & \code{+0x000} & \code{+0x080} & 未用（\code{b .}） \\
    \code{0x200} & 当前 EL，使用 \reg{SP\_ELx} & \code{el1sync} & \code{el1irq} & 内核态异常/中断 \\
    \code{0x400} & 低 EL，AArch64 & \code{el0sync} & \code{el0irq} & 系统调用、用户异常/中断 \\
    \code{0x600} & 低 EL，AArch32 & \code{+0x600} & \code{+0x680} & 未用 \\
    \bottomrule
  \end{tabular}
\end{table}

每组的第三、四个入口是 FIQ 和 SError。\file{kernel/trapasm.S} 用 \code{.balign} 保证对齐，每个入口只放一条跳转指令，真正的处理放在表外（\code{0x80} 字节只够放 32 条指令）：

\begin{asmcode}
.global alltraps
.balign 0x800
alltraps:
// Current EL with sp0
  b .
.balign 0x80
  b .
.balign 0x80
  b .
.balign 0x80
  b .

// Current EL with spx
.balign 0x80
  b el1sync
.balign 0x80
  b el1irq
.balign 0x80
  b .
.balign 0x80
  b .

// Lower EL using aarch64
.balign 0x80
  b el0sync
.balign 0x80
  b el0irq
.balign 0x80
  b .
.balign 0x80
  b .

// Lower EL using aarch32
  ...                       // 四个 b .
\end{asmcode}

只有四个入口真正被使用。其余 12 个入口是 \code{b .}，一旦进入就原地停住。这在调试初期足够了，但也有缺点：FIQ 或 SError 发生时，现象只是“毫无输出地死机”。更好的做法是让这些入口也保存现场，然后打印入口编号、\reg{ESR\_EL1} 和 \reg{ELR\_EL1} 再停机。

\subsection{保存寄存器状态}

处理完异常回到原来的代码时，所有通用寄存器必须和异常发生前一模一样，否则一个与当前代码毫不相干的中断就会悄悄改变它的行为。进入异常时，硬件只做三件事：把返回地址存入 \reg{ELR\_EL1}，把原来的 PSTATE 存入 \reg{SPSR\_EL1}，换到 \reg{SP\_EL1}。通用寄存器要由软件自己保存。

\file{trapasm.S} 为两种来源准备了两对宏。来自内核的异常用 \code{savereg}/\code{restorereg}，现场压在当前内核栈上：

\begin{asmcode}
.macro savereg
        sub sp, sp, #272              // 17 组 × 16 字节
        stp x0, x1, [sp, #16 * 0]
        stp x2, x3, [sp, #16 * 1]
        ...
        stp x28, x29, [sp, #16 * 14]
        mrs x9, elr_el1
        mrs x10, spsr_el1
        add x11, sp, #272             // 异常发生前的内核 sp
        stp x30, x9, [sp, #16 * 15]
        stp x10, x11, [sp, #16 * 16]
.endm
\end{asmcode}

来自用户态的异常用 \code{usavereg}/\code{urestorereg}，区别是最后一个字保存的是用户栈指针 \reg{SP\_EL0}，返回时也要写回：

\begin{asmcode}
.macro usavereg
        sub sp, sp, #272
        stp x0, x1, [sp, #16 * 0]
        ...
        stp x28, x29, [sp, #16 * 14]
        mrs x9, elr_el1
        mrs x10, spsr_el1
        mrs x11, sp_el0               // 用户栈指针
        stp x30, x9, [sp, #16 * 15]
        stp x10, x11, [sp, #16 * 16]
.endm

.macro urestorereg
        ldp x30, x9, [sp, #16 * 15]
        ldp x10, x11, [sp, #16 * 16]
        msr elr_el1, x9
        msr spsr_el1, x10
        msr sp_el0, x11
        ldp x0, x1, [sp, #16 * 0]
        ...
        ldp x28, x29, [sp, #16 * 14]
        add sp, sp, #272
.endm
\end{asmcode}

这 272 字节的布局与 \code{struct trapframe}（\code{x0}～\code{x30}、\code{elr}、\code{spsr}、\code{sp}）完全一致。\code{allocproc()} 把每个进程的 trapframe 放在它内核栈的最顶端，而用户态陷入时内核栈恰好是空的，所以 \code{usavereg} 压栈的位置正好就是 \code{p->trapframe}。\file{trap.c} 开头的注释画出了两种现场同时存在时一个内核栈的样子：

\begin{txtcode}
低地址   p->kstack
           |  普通内核栈空间
           |  +---------------------------+  <- 内核态中断后 sp = S - 272
           |  | el1irq savereg 临时现场   |
           |  | x0..x30, elr, spsr, sp    |
           |  +---------------------------+  <- 中断前 sp = S
           |  函数调用帧、局部变量……
           +------------------------------
           |  struct trapframe            |  <- p->trapframe（usavereg 的位置）
           |  x0..x30, elr, spsr, 用户 sp |
高地址   p->kstack + PGSIZE
\end{txtcode}

也就是说：进程在内核里执行系统调用时又被中断，会在同一个内核栈上再压一层 \code{savereg} 现场；如果中断发生在调度器中，用的则是这个 CPU 自己的调度器栈（\code{stack0}）。

有了调度器之后，“以一个进程的身份进入中断、以另一个进程的身份离开”是完全合法的：\code{eret} 依赖的 \code{elr} 和 \code{spsr} 都保存在当前进程自己的内核栈上，所以 \reg{ELR\_EL1}、\reg{SPSR\_EL1} 必须和通用寄存器一起保存和恢复。新进程第一次回到用户态走的是 \code{usertrapret}，它直接把 \code{sp} 指向 trapframe，借用 \code{el0irq} 的后半段返回：

\begin{asmcode}
el0irq:
        usavereg
        mov x0, sp
        bl userirq
trapret:
        urestorereg
        eret

.global usertrapret
usertrapret:                    // void usertrapret(struct trapframe *tf)
        mov sp, x0
        b trapret
\end{asmcode}

\subsection{设置向量表}

向量表准备好了，还要告诉处理器它在哪里：把地址写入 \reg{VBAR\_EL1}（向量基址寄存器）。这是每个核各自的系统寄存器，所以每个 CPU 都要执行一次：

\begin{ccode}
void trapinithart(void)
{
  w_vbar_el1((uint64)alltraps);
}

static inline void w_vbar_el1(uint64 x)     // kernel/aarch64.h
{
  asm volatile("msr vbar_el1, %0" : : "r" (x) );
}
\end{ccode}

\code{alltraps} 是链接时确定的高虚拟地址，所以这一步必须在 MMU 打开、跳到高地址之后进行。主核在 \code{main()} 中依次调用 \code{trapinit()}（只初始化 \code{tickslock}）和 \code{trapinithart()}；副核被唤醒后只调用 \code{trapinithart()}。

\subsection{屏蔽与解除屏蔽中断}

有些代码绝不能被中断打断。例如 \code{usavereg} 执行到一半时若又来一个中断，现场就会被覆盖。所以硬件在进入任何异常时都会自动屏蔽中断；保存好现场之后再解除屏蔽、允许嵌套，也是完全合法的。

PSTATE 中有四个屏蔽位：D（调试异常）、A（SError，也叫异步中止）、I（IRQ）、F（FIQ），用 \code{msr daifset/daifclr, \#imm} 设置或清除，立即数是一个 4 位掩码，从高到低对应 D、A、I、F。xv6 的封装在 \file{kernel/aarch64.h} 中：

\begin{ccode}
static inline void intr_on()  { asm volatile("msr daifclr, #0xf" ::: "memory"); }
static inline void intr_off() { asm volatile("msr daifset, #0xf" ::: "memory"); }

static inline int intr_get()          // IRQ 是否打开？
{
  uint64 x = daif();
  return ((x >> 6) & 0x2) == 0;       // DAIF 寄存器中 I 在 bit 7
}
\end{ccode}

注意这里用的是 \code{\#0xf}，一次改变全部四位，而不只是 I 位（\code{\#2}）。对 xv6 来说效果相同，因为它只用 IRQ；更精确的写法应只操作 I 位，以免意外打开调试异常或 SError。

\subsubsection{自旋锁与中断的嵌套}


如果一个 CPU 持有某把自旋锁时被中断，而中断处理程序又去获取同一把锁，这个 CPU 就会永远自旋下去。所以 xv6 规定：持有任何自旋锁期间，本 CPU 的中断都是关闭的。\code{acquire()} 调用 \code{push\_off()}，\code{release()} 调用 \code{pop\_off()}，二者成对嵌套计数：

\begin{ccode}
void push_off(void)
{
  int old = intr_get();
  intr_off();
  if(mycpu()->noff == 0)
    mycpu()->intena = old;          // 只记录最外层进入前的状态
  mycpu()->noff += 1;
}

void pop_off(void)
{
  struct cpu *c = mycpu();
  if(intr_get())
    panic("pop_off - interruptible");
  if(c->noff < 1)
    panic("pop_off");
  c->noff -= 1;
  if(c->noff == 0 && c->intena)     // 最外层释放、且原来是开着的，才重新打开
    intr_on();
}
\end{ccode}

\subsubsection{中断在哪里被打开}


\begin{itemize}
  \item 启动时，\file{entry.S} 通过 \code{SPSR = 0x3c5} 进入 EL1，四位全部屏蔽；
  \item 每个 CPU 第一次真正接收中断，是在 \code{scheduler()} 循环中调用 \code{intr\_on()}；
  \item 用户态陷入后，\code{usertrap()} 识别出系统调用，在调用 \code{syscall()} 之前打开中断，让长时间的系统调用也能被设备中断和定时器打断；返回用户态前再 \code{intr\_off()}，因为 \code{urestorereg} 期间不能被打断；
  \item \code{kerneltrap()} 进入时检查 \code{intr\_get()}，若中断是开着的就说明某处逻辑出错，直接 \code{panic}。
\end{itemize}

\subsection{配置中断控制器}

设备一般不直接中断 CPU，而是经过中断控制器：控制器负责打开或关闭各个中断源，并告诉 CPU 当前是谁在请求中断。BCM2837 没有 ARM 标准的 GIC，而是两个级联的模块（图~\ref{fig:irq}）：

\begin{itemize}
  \item \textbf{legacy 中断控制器}（\code{0x3f00b000}）管理 SoC 共享外设：Mini UART（IRQ 29）、DWC2 USB（IRQ 9）、Arasan SDIO（IRQ 62）、Unicam CSI1（IRQ 39）等。中断号 0～31 由 \code{ENABLE\_IRQS\_1}/\code{IRQ\_PENDING\_1} 管理，32～63 由 \code{ENABLE\_IRQS\_2}/\code{IRQ\_PENDING\_2} 管理。所有外设中断汇总成一根 “GPU IRQ”；
  \item \textbf{ARM-local 中断路由}（\code{0x40000000}）位于 CPU 簇旁，管理每个核私有的中断源：通用定时器、核间 mailbox、PMU，以及从 legacy 控制器汇总来的 GPU IRQ，最终连到四个核的 IRQ/FIQ 输入。
\end{itemize}

\begin{figure}[htbp]
  \centering
  \begin{tikzpicture}[node distance=5mm and 12mm]
    \node[box, minimum width=28mm] (uart) {Mini UART (29)};
    \node[box, minimum width=28mm, below=2mm of uart] (usb) {DWC2 USB (9)};
    \node[box, minimum width=28mm, below=2mm of usb] (sdio) {Arasan SDIO (62)};
    \node[box, minimum width=28mm, below=2mm of sdio] (csi) {Unicam CSI1 (39)};
    \node[hbox, right=of usb, yshift=-3mm, minimum height=22mm, text width=24mm] (legacy) {legacy\\中断控制器\\\code{0x3f00b000}};
    \foreach \n in {uart,usb,sdio,csi} \draw[arr] (\n.east) -- (legacy.west |- \n.east);
    \node[box, above=12mm of legacy, minimum width=26mm] (timer) {通用定时器 CNTV};
    \node[box, above=2mm of timer, minimum width=26mm] (mbox) {mailbox 0（IPI）};
    \node[hbox, right=of legacy, minimum height=40mm, text width=26mm, yshift=12mm] (local) {ARM-local\\中断路由\\\code{0x40000000}\\[2mm]\footnotesize bit 3：虚拟定时器\\bit 4：mailbox 0\\bit 8：GPU IRQ};
    \draw[arr] (legacy.east) -- node[lab, below]{GPU IRQ} (legacy.east -| local.west);
    \draw[arr] (timer.east) -- (timer.east -| local.west);
    \draw[arr] (mbox.east) -- (mbox.east -| local.west);
    \node[box, right=of local, text width=16mm] (cpus) {CPU0～3\\IRQ 输入};
    \draw[arr] (local) -- (cpus);
  \end{tikzpicture}
  \caption{BCM2837 的两级中断结构}
  \label{fig:irq}
\end{figure}

为了少改 \file{trap.c}，项目保留了原 GIC 版本的函数名 \code{gicv3init()}、\code{gicv3inithart()}、\code{gic\_iar()}、\code{gic\_eoi()}，实现换成了 \file{kernel/bcm2837.c}。全局初始化先关掉 legacy 控制器的全部中断，再只打开已经有驱动的几个：

\begin{ccode}
void gicv3init(void)
{
  irqwrite(DISABLE_IRQS_1, ~0U);
  irqwrite(DISABLE_IRQS_2, ~0U);
  irqwrite(DISABLE_BASIC_IRQS, ~0U);
  irqwrite(ENABLE_IRQS_1, (1U << UART0_IRQ) | (1U << DWC2_IRQ));
  irqwrite(ENABLE_IRQS_2, 1U << (SDIO_IRQ - 32));
}
\end{ccode}

后来加入的驱动（如摄像头）在 probe 成功后调用 \code{bcm2837\_enable\_irq(irq)} 打开对应位，按中断号小于还是大于 32 选择 \code{ENABLE\_IRQS\_1} 或 \code{\_2}。

每个核还要各自配置 ARM-local 路由：把本核的虚拟定时器接到 IRQ，并打开 mailbox 0 作为“重新调度”核间中断：

\begin{ccode}
#define CORE_TIMER_CONTROL(c) (0x40 + 4 * (c))
#define CORE_MBOX_CONTROL(c)  (0x50 + 4 * (c))
#define CORE_IRQ_SOURCE(c)    (0x60 + 4 * (c))
#define CORE_MBOX0_SET(c)     (0x80 + 0x10 * (c))   // 写 1 发出
#define CORE_MBOX0_CLR(c)     (0xc0 + 0x10 * (c))   // 写 1 清除
#define CORE_VIRTUAL_TIMER    (1 << 3)
#define CORE_MBOX0            (1 << 4)

void gicv3inithart(void)
{
  int cpu = cpuid();
  localwrite(CORE_TIMER_CONTROL(cpu), CORE_VIRTUAL_TIMER);
  localwrite(CORE_MBOX0_CLR(cpu), ~0U);     // 先清掉遗留的 mailbox 位
  localwrite(CORE_MBOX_CONTROL(cpu), 1);
}

void send_resched_ipi(int cpu)              // 让 cpu 尽快重新调度
{
  asm volatile("dsb sy" ::: "memory");      // 先让 need_resched 的写入可见
  localwrite(CORE_MBOX0_SET(cpu), 1);
}
\end{ccode}

一个外设中断要真正到达 CPU，三道开关缺一不可：设备自己的中断使能位（例如 Mini UART 的 \code{MU\_IER}）、legacy 控制器的 \code{ENABLE\_IRQS} 位，以及 CPU 的 PSTATE.I $=0$。

\subsection{通用 IRQ 处理程序}

所有 IRQ 都进入同一个入口，按来源分为两条路径：

\begin{asmcode}
el1irq:                    // 内核执行时来的 IRQ
        savereg
        mov x0, sp
        bl kernelirq
        restorereg
        eret

el0irq:                    // 用户执行时来的 IRQ
        usavereg
        mov x0, sp
        bl userirq
trapret:
        urestorereg
        eret
\end{asmcode}

两个 C 处理函数都先调用 \code{devintr()} 找出并服务中断源，再决定是否让出 CPU：

\begin{ccode}
void userirq(void)
{
  struct proc *p = myproc();
  devintr();
  if(resched_pending())      // 时间片用完，或更高优先级进程就绪（定时器或 IPI）
    yield();
  if(p->killed)
    exit(-1);
  intr_off();
}

void kernelirq()
{
  devintr();
  // 内核态也可以被抢占：只有不持任何自旋锁时中断才是开着的
  if(myproc() != 0 && myproc()->state == RUNNING && resched_pending())
    yield();
}
\end{ccode}

\code{devintr()} 先问中断控制器“是谁”，再调用对应设备的处理函数，返回 2 表示定时器、1 表示其他设备、0 表示无法识别：

\begin{ccode}
int devintr()
{
  uint32 iar = gic_iar();
  uint32 irq = gic_iar_irq(iar);
  int dev = 0;

  if(irq == UART0_IRQ){         uartintr();        dev = 1; }
  else if(irq == DWC2_IRQ){     dwc2_irq();        dev = 1; }
  else if(irq == SDIO_IRQ){     arasan_sdio_irq(); dev = 1; }
  else if(irq == CSI1_IRQ){     unicam_irq();      dev = 1; }
  else if(irq == TIMER0_IRQ){
    int logical_tick = timerintr();
    if(cpuid() == 0 && logical_tick)
      clockintr();
    sched_tick();
    dev = 2;
  } else if(irq == IPI_RESCHED_IRQ){
    resched_ipi_ack();          // need_resched 已由发送方设置
    dev = 1;
  } else if(irq == 1023){
    // 没有可处理的中断
  } else if(irq){
    printf("unexpected interrupt irq=%d\n", irq);
  }
  if(dev)
    gic_eoi(iar);
  return dev;
}
\end{ccode}

“是谁”的判断在 \code{gic\_iar()} 中。它先查本核的 ARM-local 源寄存器（定时器和 IPI 是每核私有的），再查 legacy 控制器的挂起寄存器，把两套硬件的状态翻译成一个类似 GIC 的中断号：

\begin{ccode}
uint32 gic_iar(void)
{
  int cpu = cpuid();
  uint32 local = localread(CORE_IRQ_SOURCE(cpu));
  if(local & CORE_VIRTUAL_TIMER)  return TIMER0_IRQ;        // 27
  if(local & CORE_MBOX0)          return IPI_RESCHED_IRQ;   // 伪中断号 1000
  if(irqread(IRQ_PENDING_1) & (1U << UART0_IRQ))       return UART0_IRQ;
  if(irqread(IRQ_PENDING_1) & (1U << DWC2_IRQ))        return DWC2_IRQ;
  if(irqread(IRQ_PENDING_2) & (1U << (SDIO_IRQ - 32))) return SDIO_IRQ;
  if(irqread(IRQ_PENDING_2) & (1U << (CSI1_IRQ - 32))) return CSI1_IRQ;
  return 1023;
}
\end{ccode}

BCM2837 没有 GIC 那样的 EOI 握手，\code{gic\_eoi()} 是空函数：中断要在\emph{源设备}处清除——定时器靠重新装载，UART 靠读空接收 FIFO，USB 靠写 \code{HCINT}，IPI 靠写 mailbox 清除寄存器。设备的挂起位清除之后，legacy 控制器中对应的位才会消失。多个中断可能同时挂起，\code{devintr()} 每次只处理优先级最高的一个；其余的会在 \code{eret} 之后立即再次触发中断，所以不会丢失。

\subsubsection{同步异常：系统调用路径}


同步异常走 \code{el0sync}/\code{el1sync} 入口，保存现场的方式与 IRQ 相同，只是调用的 C 函数不同。用户程序执行 \code{svc \#0} 时，系统调用号放在 \code{x7}，参数在 \code{x0}～\code{x5}：

\begin{txtcode}
user:  mov x7, #SYS_write ; svc #0
  -> VBAR_EL1 + 0x400 -> el0sync -> usavereg（现场 = trapframe）
  -> usertrap(): ESR_EL1.EC == 0x15
       -> intr_on(); syscall(): syscalls[x7]()，返回值写回 trapframe->x0
       -> 若唤醒了更高优先级进程：yield()
  -> intr_off(); urestorereg -> eret -> 回到 svc 的下一条指令
\end{txtcode}

\begin{ccode}
void usertrap(void)
{
  struct proc *p = myproc();
  uint64 esr = r_esr_el1();
  uint64 ec = (esr >> 26) & 0x3f;
  w_esr_el1(0);
  if(ec == 21){                         // SVC from AArch64
    if(p->killed) exit(-1);
    intr_on();
    syscall();
    if(resched_pending()) yield();
  } else {
    printf("usertrap(): unexpected esr=%p ec=%p pid=%d\n", esr, ec, p->pid);
    printf("            elr=%p far=%p\n", r_elr_el1(), r_far_el1());
    ...
    uvmdump(p->vm->pagetable, p->pid, p->name, "usertrap");
    p->killed = 1;
  }
  if(p->killed) exit(-1);
  intr_off();
}
\end{ccode}

硬件保存在 \reg{ELR\_EL1} 中的已经是 \code{svc} 的下一条指令地址，所以不需要像 RISC-V 那样手工给 PC 加 4。

\begin{note}
  早期版本的 \code{usertrap()} 先读出 EC 再清零 \reg{ESR\_EL1}，打印错误信息时又重新读一次寄存器，结果永远显示 \code{ec=0}。现在的代码先把完整的 ESR 存进局部变量 \code{esr}。调试信息本身的正确性，同样需要检查。
\end{note}

\subsection{初始化定时器}

xv6-aarch64 没有使用 BCM 自己的 System Timer，而是使用 ARM 架构定义的通用定时器（Generic Timer）。它的好处是：每个核都有自己的比较器，中断经 ARM-local 的 bit 3 直接送到本核，不经过 legacy 控制器；寄存器是系统寄存器，用 \code{mrs}/\code{msr} 访问，不需要 MMIO 映射。内核使用虚拟定时器那一组：

\begin{itemize}
  \item \reg{CNTFRQ\_EL0}：计数器频率（由固件设置）；
  \item \reg{CNTVCT\_EL0}：当前虚拟计数值，所有核共享，可用来测量时间（\code{clock\_us()}、\code{delay()} 都基于它）；
  \item \reg{CNTV\_TVAL\_EL0}：距离触发还剩多少个计数，写入即开始倒计时；
  \item \reg{CNTV\_CTL\_EL0}：bit 0 ENABLE、bit 1 IMASK、bit 2 ISTATUS。
\end{itemize}

\code{timerinit()} 的顺序是“关闭—装载—打开”，避免改写比较值的过程中产生一个中间中断：

\begin{ccode}
#define CNTV_CTL_ENABLE   (1<<0)
#define CNTV_CTL_IMASK    (1<<1)

void timerinit()
{
  disable_timer();
  reload_timer();
  enable_timer();
}

static void disable_timer()
{
  uint64 c = r_cntv_ctl_el0();
  c &= ~CNTV_CTL_ENABLE;
  c |= CNTV_CTL_IMASK;
  w_cntv_ctl_el0(c);
}

static void reload_timer()
{
  uint64 interval = 10000;                                   // 微秒，即 10 ms
  uint64 interval_clk = interval * (r_cntfrq_el0() / 1000000);
  w_cntv_tval_el0(interval_clk);
}

static void enable_timer()
{
  uint64 c = r_cntv_ctl_el0();
  c |= CNTV_CTL_ENABLE;
  c &= ~CNTV_CTL_IMASK;
  w_cntv_ctl_el0(c);
}
\end{ccode}

\code{timerinit()} 由每个核各自调用。执行完之后定时器已经开始倒计时，但 PSTATE.I 仍然屏蔽着 IRQ，要等进入 \code{scheduler()} 调用 \code{intr\_on()} 时 CPU 才会真正接收它。

\subsection{处理定时器中断}

定时器到期后，中断经 ARM-local 路由到本核，\code{gic\_iar()} 返回 \code{TIMER0\_IRQ}，\code{devintr()} 调用 \code{timerintr()}：

\begin{ccode}
int timerintr()
{
  int cpu = cpuid();
  int logical_tick;

  disable_timer();
  reload_timer();                       // 重新装载即清除了中断条件
  timer_divider[cpu]++;
  logical_tick = timer_divider[cpu] == 10;
  if(logical_tick)
    timer_divider[cpu] = 0;
  if(cpu == 0)
    workqueue_timer_tick();             // 推进 delayed-work 时钟
  enable_timer();
  return logical_tick;                  // 每 10 次（100 ms）返回 1
}
\end{ccode}

一次 10\,ms 的硬件中断在 \code{devintr()} 中引出三件事：

\begin{enumerate}
  \item \textbf{系统时间}：每 10 次中断构成一个 xv6 逻辑 tick。只有 CPU0 调用 \code{clockintr()} 把全局 \code{ticks} 加一并 \code{wakeup(\&ticks)}，避免四个核把时间推进四倍。\code{sleep()}、\code{uptime()} 因此保持 xv6 原来 100\,ms 一个 tick 的语义：
\begin{ccode}
void clockintr()
{
  acquire(&tickslock);
  ticks++;
  wakeup(&ticks);
  release(&tickslock);
  net_udp_timeout_tick();       // 检查带超时的 UDP 接收是否到期
}
\end{ccode}
  \item \textbf{下半部时钟}：CPU0 推进 workqueue 的 delayed-work 时钟，TCP 重传、Wi-Fi 连接超时等到期的工作被放进各自的工作队列，由内核线程执行，而不是在中断里直接运行；
  \item \textbf{调度节拍}：每次中断都调用 \code{sched\_tick()}，扣减当前进程的时间片、累计实时任务的运行时间，必要时设置 \code{need\_resched}（第~\ref{chap:process} 章）。随后 \code{userirq()}/\code{kernelirq()} 检查 \code{resched\_pending()}，决定是否 \code{yield()}。
\end{enumerate}

\code{reload\_timer()} 写的是相对倒计时，所以下一次中断发生在“本次处理程序重新装载的时刻 $+$ 10\,ms”，处理程序的延迟会累积成时钟漂移。若要更稳定的时基，可以改写绝对比较值 \reg{CNTV\_CVAL\_EL0}，每次在上一个截止时间上加固定间隔。

\subsection{小结}

把本节的代码按 \code{main()} 中的调用顺序串起来：

\begin{ccode}
trapinit();        // 初始化 tickslock
trapinithart();    // VBAR_EL1 = alltraps                   （每个核）
gicv3init();       // legacy 控制器：全关，再打开 UART/USB/SDIO
gicv3inithart();   // ARM-local：本核虚拟定时器 + mailbox 0  （每个核）
timerinit();       // 10 ms 虚拟定时器开始倒计时            （每个核）
...
scheduler();       // intr_on()：第一次真正接收中断
\end{ccode}

顺序不能随意交换：先安装向量表，保证中断来了有入口；再配置控制器和路由；再启动定时器；最后才解除 PSTATE 的屏蔽。要收到第一个定时器中断，下面这些条件必须同时满足：

\begin{itemize}
  \item 处于 EL1h，且 EL2 已允许 EL1 访问计数器和定时器（\reg{CNTHCTL\_EL2}）；
  \item \reg{VBAR\_EL1} 指向 2\,KiB 对齐的 \code{alltraps}；
  \item \code{CORE\_TIMER\_CONTROL(cpu)} 打开了虚拟定时器路由；
  \item \reg{CNTV\_TVAL\_EL0} 装入了非零值，\reg{CNTV\_CTL\_EL0} 的 ENABLE $=1$、IMASK $=0$；
  \item PSTATE.I $=0$，且当前内核栈有效，有空间保存 272 字节的现场。
\end{itemize}

调试中断问题时，可以在启动阶段打印 \reg{CurrentEL}、\reg{VBAR\_EL1}、\reg{DAIF}、\reg{CNTFRQ\_EL0} 和 \reg{CNTV\_CTL\_EL0}；对未知中断打印 \code{IRQ\_PENDING\_1/2} 和 \code{CORE\_IRQ\_SOURCE}。但不要在定时器中断里持续 \code{printf}：串口远比中断慢，会改变调度时序，甚至带来新的锁问题。


\setcounter{secnumdepth}{2}
